% Chapter 2, section 4: Pop messages
% Write in plain LaTeX. Code of a guide step: \guide{<step>}; commands: \begin{shell}...\end{shell};
% a screenshot: \screenshot{file}{caption}; a gap in the draft: \placeholder{...}.
% See book/latex/queen.sty for the full list.

\section{Pop messages}

There are two ways to read. \texttt{consume()} is a loop: it takes messages, calls your handler,
acks them when the handler returns, and goes on. \texttt{pop()} is one call: it hands you the
messages and leaves the ack to you.

\why{most workers are loops, and \texttt{consume()} writes the loop for you, with the error handling
done right. \texttt{pop()} is for the cases a loop does not fit: a scheduled job that reads once, a
script, or code that wants to decide itself when to ack.}

\queenfig{consume-loop}{What consume() does for you.}

\subsection{Queue mode and consumer groups}

Without a consumer group, each message goes to one worker, and the broker sizes every pop for you.

\guide{consume}

With a consumer group, every group gets every message and the workers of one group share them. A
new group starts at the tail, so this one asks for everything already in the queue. Here every
option is set by hand, and the handler acks.

\guide{consume-group}

Ack \texttt{failed} and the message comes back, up to the queue's \texttt{retryLimit}; after the last retry it
goes to the dead-letter queue.

\why{one queue often feeds several services: billing, emails and analytics all want every order.
Each service is a consumer group, reading at its own pace, and none of them can slow the others
down. Inside a group, the workers compete, so you add workers to go faster.}

\queenfig{queue-mode-vs-groups}{Queue mode shares the messages; every consumer group gets all of them.}

\subsection{One message at a time, or a whole batch}

With \texttt{.each()}, the handler gets one message at a time. Without it, the handler gets the
whole batch as an array.

\guide{consume-batch}

What the client does with your handler's result (\texttt{autoAck}, on by default):

\begin{itemize}
\item \textbf{It returns:} the message (or the whole batch) is acked.
\item \textbf{It throws:} the message (or the whole batch) is nacked as \texttt{failed}, which
spends one retry, and the loop goes on with the next messages.
\item \textbf{It throws under \texttt{.each()}:} the rest of that partition's messages in the same
batch are skipped, because they will be delivered again after the failed one, in order. Messages of
other partitions in the batch are handled as usual.
\item \textbf{With \texttt{.autoAck(false)}:} the handler acks itself (in a transaction, for
example). A handler that throws is still nacked, so its partition is not stuck until the lease
runs out.
\end{itemize}

\why{a batch is cheaper (one ack for a hundred messages, one database write for a hundred rows),
and one at a time is simpler and fails one message instead of a hundred. Pick per worker.}

If you add \texttt{.onSuccess()} or \texttt{.onError()} to a consumer, the client stops acking and
nacking for you: you decide in those callbacks what happens to each message.

\guide{consume-callbacks}

\why{some failures should be retried and some never will succeed (an email address that does not
exist). The callbacks let you send the first kind back with \texttt{failed} and the second kind
straight to the dead-letter queue with \texttt{dlq}.}

\subsection{How many messages a pop takes: autopilot}

By default the broker decides how many messages each pop gets (\texttt{batch}, from 100 to 1,000,
based on how fast the group acks) and from how many partitions (\texttt{partitions}, the ones that
have work, shared among the waiting workers). A value you set yourself is never changed.
\texttt{partitions(N)} lets one pop take messages from up to \texttt{N} partitions (1 to 64), with
\texttt{batch} as the total. \texttt{.autopilot(false)} goes back to one message from one partition
per pop.

\why{the right batch size depends on your handler and your load, and it changes during the day.
The broker sees both, so it can size every pop better than a number fixed in the code.}

\subsection{Leases and long work}

A pop leases its messages to the worker for the queue's \texttt{leaseTime}. Work that can take
longer renews the lease while it runs: \texttt{.renewLease(true, 10000)} renews every 10 seconds,
and \texttt{queen.renew(message)} renews once. In PHP, the renewal happens between messages, so a
single long handler calls the renew callback it is given. A lease that has already
expired cannot be renewed: its messages go to the next worker, marked as one delivery attempt
more.

\why{the lease is what protects order and stops two workers from doing the same work. A short lease
gets a crashed worker's messages moving again quickly; renewing lets a slow but healthy worker keep
them.}

\subsection{Where a new group starts}

A group's starting point is decided on its first pop and then stored: \texttt{new} (the default)
reads only what is pushed after that first pop, \texttt{all} reads everything still kept, and
\texttt{subscriptionFrom} reads from a point in time. A group that already exists moves with a seek,
to a time or to the end.

\queenfig{subscription-start}{Where a new group starts, fixed at its first pop.}

\guide{seek}

\why{a new service usually wants only new events, a new projection wants the whole history, and a
bug fix wants to replay the last day. A seek to the end is how a group skips a backlog it no longer
needs.}

\subsection{Only the newest message: conflation}

A group with \texttt{.conflation()} gets only the newest message of each partition, and its cursor
moves past the older ones.

\guide{conflation}

\why{some queues carry commands where only the latest one matters: ``recompute the price of A-1''.
A worker that is behind should do the work once, with the newest input, instead of three times with
inputs that are already old.} Conflation is decided by the group's first pop, needs a consumer
group, and cannot be combined with \texttt{commitOnDelivery}.

\subsection{Faster, with no ack: commitOnDelivery}

A pop with \texttt{.commitOnDelivery()} moves the group's cursor as the messages are handed out.
There is no lease and nothing to ack.

\guide{commit-on-delivery}

\why{a live dashboard or a metrics chart can lose a reading now and then, and would rather not pay
for leases and acks. This is at-most-once delivery: a worker that crashes after the pop loses that
batch. It works with \texttt{pop()} only; \texttt{consume()} refuses it.}

\subsection{Running a consumer for real}

\guide{consume-run}

\begin{itemize}
\item \texttt{concurrency(n)} runs \texttt{n} handlers at once, each on its own partitions.
\item \texttt{limit(n)} stops after \texttt{n} messages, and \texttt{idleMillis(ms)} stops after
that long with nothing to do.
\item \texttt{wait(true)} (the default) waits for messages instead of returning empty, up to
\texttt{timeoutMillis(ms)} (30 seconds by default, at most 60).
\item A \texttt{signal} stops the loop at once. Messages it had taken but not started go back with a
\texttt{retry} ack, which spends no retry, so a deploy does not leave partitions waiting for their
leases to run out.
\end{itemize}

\why{a worker runs for weeks, through deploys and quiet nights. These options decide how much it
does at once, when it stops, and how it gives its work back when it is told to.}

\subsection{Read it again}

Consuming does not delete: the queue is a log, kept for its retention. A new group can start from a
point in time and read the same orders again, while the groups that already read them stay where
they are.

\guide{replay}

Or read a partition by offset with \texttt{fetch}: no group, no lease, no ack. It suits an audit,
rebuilding one entity, or a sink that keeps its own position.

\guide{fetch}

\why{history is useful long after the first reader is done with it: to build a new view, to audit
what happened, or to rebuild one customer's state from its events.}

On the site, with the code in all six SDK languages: \href{https://queenmq.com/writing-applications/pop/}{queenmq.com/writing-applications/pop}. More there: \href{https://queenmq.com/writing-applications/clients/}{reads by offset in every SDK}.
